\documentclass[conference]{IEEEtran}

% ── Packages ──────────────────────────────────────────────────────────────────
\usepackage{cite}
\usepackage{amsmath,amssymb,amsfonts}
\usepackage{algorithmic}
\usepackage{graphicx}
\usepackage{textcomp}
\usepackage{xcolor}
\usepackage{booktabs}
\usepackage{multirow}
\usepackage{array}
\usepackage{url}
\usepackage{float}
\def\BibTeX{{\rm B\kern-.05em{\sc i\kern-.025em b}\kern-.08em
    T\kern-.1667em\lower.7ex\hbox{E}\kern-.125emX}}

% ──────────────────────────────────────────────────────────────────────────────
\begin{document}

\title{Machine Learning-Based Prediction of Machining Process State and\\
Spindle Output Power from CNC Milling Sensor Data}

\author{
\IEEEauthorblockN{Zil Muarij (456149)\IEEEauthorrefmark{1},
Ayman Arshad (479690)\IEEEauthorrefmark{1},
Khwaja Qais (475156)\IEEEauthorrefmark{1},
Abdullah Naeem (465448)\IEEEauthorrefmark{1}}
\IEEEauthorblockA{\IEEEauthorrefmark{1}Department of Mechanical Engineering,\\
College of Electrical \& Mechanical Engineering (CEME),\\
National University of Sciences and Technology (NUST), Rawalpindi, Pakistan}
}

\maketitle

% ── ABSTRACT ──────────────────────────────────────────────────────────────────
\begin{abstract}
This paper presents a machine learning framework for predicting machining process states
and spindle output power from CNC milling sensor data. The dataset comprises 1,448 samples
(post-cleaning) with 48 features including axis velocities, currents, voltages, and feed
rates recorded during multi-process milling operations. Three classification models —
Logistic Regression, Support Vector Machine (SVM), and Random Forest — are trained for
dual-target prediction of machining process category and machine speed state. Additionally,
three regression models — Linear Regression, Random Forest Regressor, and Support Vector
Regressor (SVR) — predict spindle output power. Random Forest achieves the highest
classification accuracy of 65.6\% (machining process) and 100\% (machine state), and
the best regression performance with $R^2 = 0.9999$ and relative RMSE of 0.37\%.
Statistical significance tests (Wilcoxon, $p = 0.00195$) confirm Random Forest's
superiority. These results demonstrate that ensemble tree methods are highly suitable for
real-time CNC process monitoring applications.
\end{abstract}

\begin{IEEEkeywords}
CNC milling, machining process classification, spindle power prediction, Random Forest,
Support Vector Machine, machine learning, predictive manufacturing.
\end{IEEEkeywords}

% ── I. INTRODUCTION ───────────────────────────────────────────────────────────
\section{Introduction}

Modern computer numerical control (CNC) machining centers generate vast streams of
real-time sensor data covering axis positions, velocities, motor currents, voltages,
and feed rates. Harnessing this data through machine learning (ML) offers significant
opportunities for process monitoring, quality assurance, and predictive maintenance
— core pillars of Industry 4.0 manufacturing \cite{ref1}.

Manual identification of machining states (e.g., drilling, milling, idle) is error-prone
and impractical at production scale. Automated classification of process states enables
real-time anomaly detection, tool wear monitoring, and adaptive control. Simultaneously,
accurate prediction of spindle output power supports energy efficiency optimization and
avoids overloading drive systems \cite{ref2}.

Prior work has demonstrated the utility of tree-based ensembles \cite{ref3} and SVMs
\cite{ref4} for machining process classification. However, few studies simultaneously
address multi-output classification and regression on the same sensor stream with
rigorous statistical validation. This paper addresses this gap using the UC Irvine
CNC Milling dataset \cite{ref5}.

The remainder of this paper is organized as follows. Section~\ref{sec:dataset} describes
the dataset and experimental setup. Section~\ref{sec:methodology} details preprocessing
and model formulations. Section~\ref{sec:results} presents and compares all results.
Section~\ref{sec:discussion} interprets findings from an engineering perspective.
Section~\ref{sec:conclusion} concludes with future directions.

% ── II. DATASET & EXPERIMENTAL SETUP ──────────────────────────────────────────
\section{Dataset \& Experimental Setup}
\label{sec:dataset}

\subsection{Dataset Source}

The dataset used is the \textit{CNC Mill Tool Wear} dataset (experiment\_03), sourced
from the UC Irvine Machine Learning Repository \cite{ref5}. It contains time-series
sensor recordings from a CNC milling machine performing multiple distinct machining
operations, logged at fixed sampling intervals.

\subsection{Feature Description}

The raw dataset contains sensor readings from three linear axes (X1, Y1, Z1) and the
spindle (S1), covering actual/command positions, velocities, accelerations, output
currents, voltages, DC bus voltages, output power, and current feedback signals. Two
derived features were engineered (Section~\ref{sec:fe}), yielding 48 input features
after preprocessing. Table~\ref{tab:features} summarises the key feature groups.

\begin{table}[h]
\caption{Feature Groups in the CNC Milling Dataset}
\label{tab:features}
\centering
\begin{tabular}{lll}
\toprule
\textbf{Group} & \textbf{Examples} & \textbf{Type} \\
\midrule
Axis Position   & X1\_ActualPosition, Y1\_CommandPosition  & Continuous \\
Axis Velocity   & X1\_ActualVelocity, Z1\_CommandVelocity  & Continuous \\
Axis Current    & X1\_OutputCurrent, Y1\_CurrentFeedback   & Continuous \\
Axis Voltage    & X1\_OutputVoltage, Y1\_DCBusVoltage      & Continuous \\
Axis Power      & X1\_OutputPower, Y1\_OutputPower         & Continuous \\
Spindle         & S1\_OutputCurrent, S1\_OutputVoltage     & Continuous \\
Process Info    & M1\_CURRENT\_FEEDRATE, NT\_PROGRAM\_NUMBER & Mixed \\
Engineered      & Resultant\_Velocity, Spindle\_Power\_Calc & Continuous \\
\midrule
\textbf{Targets} & Machining\_Process (9 classes) & Categorical \\
                 & Machine\_State (3 classes)      & Categorical \\
                 & S1\_OutputPower                 & Continuous \\
\bottomrule
\end{tabular}
\end{table}

\subsection{Exploratory Data Analysis}

Fig.~\ref{fig:histograms} shows the distribution histograms of all 48 features. Most
axis position and velocity features exhibit bimodal or multimodal distributions,
reflecting the distinct operating regimes of different machining processes. Spindle
features (S1) show narrow, peaked distributions indicating stable spindle operation,
while feed rate and program number features display discrete distributions consistent
with their categorical/stepped nature.

\begin{figure}[h]
    \centering
    \includegraphics[width=0.95\columnwidth]{Screenshot 2026-05-03 230347.png}
    \caption{Feature distribution histograms for all 48 input variables. Bimodal
             distributions in axis features reflect multiple distinct machining states.}
    \label{fig:histograms}
\end{figure}

\subsection{Train/Validation/Test Split}

After outlier removal, 1,448 samples remained. A stratified 70/15/15 split was applied:

\begin{itemize}
    \item \textbf{Training set:} 1,013 samples (70\%)
    \item \textbf{Validation set:} 217 samples (15\%)
    \item \textbf{Test set:} 218 samples (15\%)
\end{itemize}

Stratification was performed on the \textit{Machining\_Process} label to ensure class
balance across all splits. All models used identical splits with \texttt{random\_state=42}
for full reproducibility.

% ── III. METHODOLOGY ──────────────────────────────────────────────────────────
\section{Methodology}
\label{sec:methodology}

\subsection{Preprocessing}

\textbf{Missing Value Imputation:} Missing values were imputed using column means —
appropriate for sensor data where missingness is random due to transmission gaps rather
than systematic failure.

\textbf{Outlier Removal:} Z-score based outlier detection was applied to all numeric
features. Rows where more than 5\% of features exceeded $|z| > 3$ were removed:
\begin{equation}
    z_i = \frac{x_i - \mu}{\sigma}, \quad \text{remove if } \frac{1}{p}\sum_{j=1}^{p}
    \mathbf{1}[|z_{ij}|>3] > 0.05
    \label{eq:zscore}
\end{equation}
This retained 1,448 of the original samples, removing approximately 10\% as statistical
outliers. Samples at true operating extremes were retained to preserve physical validity.

\textbf{Normalization:} All input features were standardized using zero-mean,
unit-variance scaling:
\begin{equation}
    x' = \frac{x - \mu_{\text{train}}}{\sigma_{\text{train}}}
    \label{eq:scaling}
\end{equation}
Scaling parameters were computed on the training set only, then applied to validation
and test sets to prevent data leakage.

\subsection{Feature Engineering}
\label{sec:fe}

Two physically meaningful features were derived:

\textbf{Feature 1 — Resultant Velocity:} The Euclidean norm of the three axis velocities
represents the true magnitude of tool motion in 3D space:
\begin{equation}
    v_{\text{res}} = \sqrt{v_{X}^2 + v_{Y}^2 + v_{Z}^2}
    \label{eq:vres}
\end{equation}
This feature was subsequently used to define the \textit{Machine\_State} target via
tertile quantisation (Low/Medium/High speed).

\textbf{Feature 2 — Spindle Power (Calculated):} Electrical power delivered to the
spindle motor was computed from measured current and voltage:
\begin{equation}
    P_{\text{spindle}} = I_{S1} \times V_{S1}
    \label{eq:pspindle}
\end{equation}
This provides an independent estimate of spindle load, distinct from the directly
logged \texttt{S1\_OutputPower}, and serves as a physical cross-check feature.

\subsection{Classification Models}

\textbf{Logistic Regression:} A multi-class linear classifier using the softmax
function. For class $k$ with weight vector $\mathbf{w}_k$:
\begin{equation}
    P(y=k \mid \mathbf{x}) = \frac{e^{\mathbf{w}_k^T \mathbf{x}}}
    {\sum_{j=1}^{K} e^{\mathbf{w}_j^T \mathbf{x}}}
    \label{eq:logreg}
\end{equation}
Regularization parameter $C=10$ was selected via 5-fold cross-validation grid search.
Implemented as \texttt{MultiOutputClassifier} to handle dual targets simultaneously.

\textbf{Support Vector Machine (SVM):} A kernel-based classifier that maximizes the
margin between classes. The decision boundary solves:
\begin{equation}
    \min_{\mathbf{w},b,\xi} \frac{1}{2}\|\mathbf{w}\|^2 + C\sum_{i=1}^{n}\xi_i
    \label{eq:svm}
\end{equation}
subject to $y_i(\mathbf{w}^T\phi(\mathbf{x}_i)+b) \geq 1-\xi_i$. A linear kernel
with $C=10$ was selected via grid search.

\textbf{Random Forest:} An ensemble of $T$ decision trees trained on bootstrap samples,
with predictions aggregated by majority vote:
\begin{equation}
    \hat{y} = \text{mode}\{h_t(\mathbf{x})\}_{t=1}^{T}
    \label{eq:rf}
\end{equation}
Optimal hyperparameters: $T=200$ trees, \texttt{max\_depth}$=\infty$ (fully grown),
selected via 5-fold cross-validation. Class weights were balanced to handle class
imbalance in the machining process labels.

\subsection{Regression Models}

\textbf{Linear Regression:} Ordinary least squares estimation of spindle output power:
\begin{equation}
    \hat{y} = \mathbf{w}^T\mathbf{x} + b, \quad
    \mathbf{w}^* = \arg\min_{\mathbf{w}} \|\mathbf{X}\mathbf{w} - \mathbf{y}\|^2
    \label{eq:linreg}
\end{equation}

\textbf{Random Forest Regressor:} Ensemble of regression trees with predictions
averaged across $T=100$ trees:
\begin{equation}
    \hat{y} = \frac{1}{T}\sum_{t=1}^{T} h_t(\mathbf{x})
    \label{eq:rfr}
\end{equation}

\textbf{Support Vector Regressor (SVR):} Fits a function within an $\varepsilon$-tube
around the training data:
\begin{equation}
    \min_{\mathbf{w}} \frac{1}{2}\|\mathbf{w}\|^2 + C\sum_{i=1}^{n}(\xi_i + \xi_i^*)
    \label{eq:svr}
\end{equation}
subject to $|y_i - \mathbf{w}^T\phi(\mathbf{x}_i) - b| \leq \varepsilon + \xi_i^{(*)}$.

% ── IV. RESULTS ───────────────────────────────────────────────────────────────
\section{Results}
\label{sec:results}

\subsection{Classification Performance}

Table~\ref{tab:clf} presents the classification results on the held-out test set for
all three models across both targets. Random Forest substantially outperforms the
linear methods on Target 1 (Machining Process), achieving 65.6\% accuracy versus
26.1\% for Logistic Regression and 26.6\% for SVM. All models perform comparably on
the engineered Target 2 (Machine State), achieving $\geq$96.8\% accuracy — a result
expected given its deterministic construction from resultant velocity.

\begin{table}[h]
\caption{Classification Results on Test Set (218 Samples)}
\label{tab:clf}
\centering
\begin{tabular}{lcccc}
\toprule
\multirow{2}{*}{\textbf{Model}} &
\multicolumn{2}{c}{\textbf{Target 1 (Process)}} &
\multicolumn{2}{c}{\textbf{Target 2 (State)}} \\
\cmidrule(lr){2-3}\cmidrule(lr){4-5}
 & Acc. & F1 (W) & Acc. & F1 (W) \\
\midrule
Logistic Regression & 0.261 & 0.26 & 0.972 & 0.97 \\
SVM (Linear)        & 0.266 & 0.27 & 0.968 & 0.97 \\
Random Forest       & \textbf{0.656} & \textbf{0.66} & \textbf{1.000} & \textbf{1.00} \\
\bottomrule
\end{tabular}
\end{table}

Detailed per-class metrics for Random Forest (Target 1) are given in
Table~\ref{tab:rf_detail}, confirming strong performance across most machining classes.

\begin{table}[h]
\caption{Random Forest Per-Class Metrics — Machining Process (Target 1)}
\label{tab:rf_detail}
\centering
\begin{tabular}{ccccc}
\toprule
\textbf{Class} & \textbf{Prec.} & \textbf{Rec.} & \textbf{F1} & \textbf{Support} \\
\midrule
0 & 0.53 & 0.56 & 0.55 & 16 \\
1 & 0.76 & 0.62 & 0.68 & 26 \\
2 & 0.59 & 0.55 & 0.57 & 31 \\
3 & 0.74 & 0.72 & 0.73 & 47 \\
4 & 0.80 & 0.57 & 0.67 & 14 \\
5 & 0.74 & 0.88 & 0.80 & 16 \\
6 & 0.45 & 0.54 & 0.49 & 24 \\
7 & 1.00 & 0.50 & 0.67 & 10 \\
8 & 0.64 & 0.79 & 0.71 & 34 \\
\midrule
\textbf{Weighted Avg} & \textbf{0.67} & \textbf{0.66} & \textbf{0.66} & \textbf{218} \\
\bottomrule
\end{tabular}
\end{table}

\subsection{Regression Performance}

Table~\ref{tab:reg} compares the three regression models. Linear Regression and Random
Forest Regressor both achieve near-perfect $R^2 > 0.9998$, while SVR fails significantly
($R^2 = 0.164$) due to its sensitivity to the default RBF kernel parameters and the
narrow power range in the data.

\begin{table}[h]
\caption{Regression Results — Spindle Output Power Prediction}
\label{tab:reg}
\centering
\begin{tabular}{lcccc}
\toprule
\textbf{Model} & \textbf{RMSE} & \textbf{MAE} & $\mathbf{R^2}$ & \textbf{Adj.}~$\mathbf{R^2}$ \\
\midrule
Linear Regression    & 8.89e-4 & 6.24e-4 & 0.9999 & 0.9998 \\
Random Forest        & \textbf{8.24e-4} & \textbf{5.52e-4} & \textbf{0.9999} & \textbf{0.9999} \\
SVR                  & 7.37e-2 & 6.96e-2 & 0.164  & $-$0.073 \\
\bottomrule
\end{tabular}
\end{table}

\subsection{Learning Curves}

Figs.~\ref{fig:lc_logistic}, \ref{fig:lc_svm}, and \ref{fig:lc_rf} show the learning
curves for all three classification models. Logistic Regression and SVM both exhibit a
large persistent gap between training and validation accuracy, indicative of underfitting
on this non-linear, high-dimensional dataset. Random Forest achieves perfect training
accuracy (1.00) with a converging validation curve reaching approximately 0.80 at full
training size, suggesting mild overfitting that is expected to reduce with further data.

\begin{figure}[h]
    \centering
    \includegraphics[width=0.95\columnwidth]{Screenshot 2026-05-03 230956.png}
    \caption{Learning curve — Logistic Regression. Large train/validation gap indicates
             underfitting due to linear model insufficiency on non-linear CNC data.}
    \label{fig:lc_logistic}
\end{figure}

\begin{figure}[h]
    \centering
    \includegraphics[width=0.95\columnwidth]{Screenshot 2026-05-03 231033.png}
    \caption{Learning curve — SVM (Linear kernel). Decreasing training accuracy with
             increasing data size confirms the model is underfitting the feature space.}
    \label{fig:lc_svm}
\end{figure}

\begin{figure}[h]
    \centering
    \includegraphics[width=0.95\columnwidth]{Screenshot 2026-05-03 231101.png}
    \caption{Learning curve — Random Forest. Training accuracy = 1.00 throughout;
             validation accuracy converges toward 0.80, suggesting mild overfitting
             that decreases with additional training data.}
    \label{fig:lc_rf}
\end{figure}

\subsection{Residual Analysis}

Fig.~\ref{fig:residual} shows the residual plot for the best regression model (Random
Forest). Residuals are randomly distributed around zero with no systematic trend,
confirming homoscedasticity and absence of model bias. Fig.~\ref{fig:avp} shows near-perfect
alignment of predicted versus actual spindle output power values.

\begin{figure}[h]
    \centering
    \includegraphics[width=0.95\columnwidth]{Screenshot 2026-05-03 230926.png}
    \caption{Residual plot for Random Forest Regressor. Residuals are randomly scattered
             around zero, indicating no systematic prediction bias.}
    \label{fig:residual}
\end{figure}

\begin{figure}[h]
    \centering
    \includegraphics[width=0.95\columnwidth]{Screenshot 2026-05-03 230933.png}
    \caption{Actual vs.\ predicted spindle output power — Random Forest Regressor.
             Points align closely along the diagonal, confirming $R^2 = 0.9999$.}
    \label{fig:avp}
\end{figure}

\subsection{Correlation Heatmap}

Fig.~\ref{fig:heatmap} shows the feature correlation heatmap. Strong intra-axis
correlations are visible (e.g., output power, current, and voltage within the same
axis), while cross-axis correlations are weaker — justifying retention of all sensor
channels as independent information sources.

\begin{figure}[h]
    \centering
    \includegraphics[width=0.95\columnwidth]{Screenshot 2026-05-03 230355.png}
    \caption{Pearson correlation heatmap of all 48 features. Diagonal blocks indicate
             strong intra-axis correlations; off-diagonal blocks are largely uncorrelated.}
    \label{fig:heatmap}
\end{figure}

\subsection{Statistical Significance Tests}

To determine whether Random Forest's superiority is statistically significant, paired
10-fold cross-validation scores were compared using both a paired $t$-test and the
Wilcoxon signed-rank test (Table~\ref{tab:stats}).

\begin{table}[h]
\caption{Statistical Significance: RF vs.\ Baseline Models}
\label{tab:stats}
\centering
\begin{tabular}{lccc}
\toprule
\textbf{Comparison} & \textbf{RF Mean CV} & \textbf{$p$ (t-test)} & \textbf{$p$ (Wilcoxon)} \\
\midrule
RF vs.\ SVM      & 0.817 vs.\ 0.626 & $4.97 \times 10^{-8}$ & 0.00195 \\
RF vs.\ Logistic & 0.817 vs.\ 0.607 & $2.28 \times 10^{-9}$ & 0.00195 \\
\bottomrule
\end{tabular}
\end{table}

Both tests yield $p \ll 0.05$, confirming that Random Forest is \textit{statistically
significantly} superior to both SVM and Logistic Regression — not merely numerically
better by chance.

% ── V. DISCUSSION ─────────────────────────────────────────────────────────────
\section{Discussion}
\label{sec:discussion}

\subsection{Engineering Interpretation}

\textbf{Classification:} The low accuracy of Logistic Regression (26.1\%) and SVM
(26.6\%) on the 9-class machining process target reflects the fundamentally non-linear
decision boundaries required. CNC machining states are separated by complex combinations
of axis dynamics that linear and kernel-linear models cannot adequately capture.
Random Forest's 65.6\% accuracy is a meaningful result for a 9-class problem with
highly overlapping sensor signatures between similar operations (e.g., roughing vs.
finishing passes on the same axis).

\textbf{Regression:} The Random Forest regressor achieves a relative RMSE of only
0.37\%, qualifying as \textit{excellent} for engineering deployment. In physical units,
the RMSE of $8.24 \times 10^{-4}$ kW (0.824 W) is negligible compared to typical
spindle powers in the 140--220 W range observed in this dataset — well within
measurement noise. Linear regression performs almost identically ($R^2 = 0.9998$),
indicating that the relationship between input features and spindle output power is
predominantly linear. SVR's poor performance ($R^2 = 0.164$) is attributed to the
default RBF kernel being poorly matched to the linear power relationship and the
very narrow output range ($\sim$0 to 0.22 W).

\subsection{Feature Importance}

For Random Forest classification, the top predictive features are
\texttt{S1\_CommandPosition} (0.063), \texttt{S1\_ActualPosition} (0.061), and
\texttt{Y1\_CurrentFeedback} (0.043), indicating that spindle position and axis
current feedback are the primary discriminators of machining state. This aligns with
engineering intuition: the spindle position encodes which part of the machining program
is being executed, while current feedback reflects the actual cutting load.

\subsection{Limitations and Failure Analysis}

\begin{itemize}
    \item \textbf{Class imbalance:} The 9-class machining process target has unequal
          support (10--47 samples per class), contributing to lower recall on
          minority classes (e.g., class 0: recall = 0.56, class 7: recall = 0.50).
    \item \textbf{RF overfitting:} The training/validation gap in the RF learning curve
          (Fig.~\ref{fig:lc_rf}) suggests overfitting. Applying SMOTE for minority
          classes and tuning \texttt{min\_samples\_leaf} could reduce this.
    \item \textbf{SVR failure:} The SVR regressor requires explicit kernel and
          hyperparameter tuning (RBF $\gamma$, $\varepsilon$) to be competitive.
          Without this, it cannot match linear models on near-linear targets.
    \item \textbf{Dataset scope:} Results are specific to the machining conditions
          in this experiment. Generalization to different materials, tools, or machines
          requires transfer learning or domain adaptation.
\end{itemize}

% ── VI. CONCLUSION ────────────────────────────────────────────────────────────
\section{Conclusion}
\label{sec:conclusion}

This paper demonstrated a complete machine learning pipeline for CNC milling process
monitoring using multi-axis sensor data. Among three classification models, Random
Forest achieved the highest accuracy of 65.6\% for 9-class machining process
identification and 100\% for 3-class machine speed state prediction. For spindle
output power regression, Random Forest and Linear Regression both achieved
$R^2 > 0.9998$ with a relative RMSE of 0.37\% — well within engineering acceptance
limits. Statistical significance tests (Wilcoxon $p = 0.00195$) confirmed that Random
Forest's advantage over SVM and Logistic Regression is not due to chance. Based on
these findings, \textbf{Random Forest is recommended for deployment} in real-time
CNC process monitoring systems.

Future work should investigate: (1) deep learning approaches (LSTM, 1D-CNN) for
temporal sequence modeling of sensor streams; (2) SMOTE-based oversampling for
minority machining classes; (3) online/incremental learning for adaptation to new
tools and materials; and (4) model deployment on edge hardware for real-time
in-process monitoring.

% ── REFERENCES ────────────────────────────────────────────────────────────────
\begin{thebibliography}{00}

\bibitem{ref1}
J. Lee, B. Bagheri, and H. A. Kao, ``A cyber-physical systems architecture for industry
4.0-based manufacturing systems,'' \textit{Manufacturing Letters}, vol.~3, pp.~18--23,
2015.

\bibitem{ref2}
Y. Altintas, \textit{Manufacturing Automation: Metal Cutting Mechanics, Machine Tool
Vibrations, and CNC Design}, 2nd ed. Cambridge: Cambridge University Press, 2012.

\bibitem{ref3}
L. Breiman, ``Random forests,'' \textit{Machine Learning}, vol.~45, no.~1,
pp.~5--32, 2001.

\bibitem{ref4}
C. Cortes and V. Vapnik, ``Support-vector networks,''
\textit{Machine Learning}, vol.~20, no.~3, pp.~273--297, 1995.

\bibitem{ref5}
M. Dua and C. Graff, ``CNC Mill Tool Wear Dataset,''
UCI Machine Learning Repository, 2019. [Online]. Available:
\url{https://archive.ics.uci.edu/ml/datasets/CNC+Mill+Tool+Wear}

\bibitem{ref6}
T. Chen and C. Guestrin, ``XGBoost: A scalable tree boosting system,''
in \textit{Proc. 22nd ACM SIGKDD Conf. Knowledge Discovery and Data Mining},
pp.~785--794, 2016.

\bibitem{ref7}
F. Pedregosa \textit{et al.}, ``Scikit-learn: Machine learning in Python,''
\textit{Journal of Machine Learning Research}, vol.~12, pp.~2825--2830, 2011.

\bibitem{ref8}
G. Georgoulas, T. Loutas, C. D. Stylios, and V. Kostopoulos, ``Bearing fault detection
based on hybrid ensemble detector and empirical mode decomposition,''
\textit{Mechanical Systems and Signal Processing}, vol.~41, pp.~510--525, 2013.

\bibitem{ref9}
A. Agogino and K. Goebel, ``CMAPSS Jet Engine Simulated Data,''
NASA Ames Prognostics Data Repository, 2007. [Online]. Available:
\url{https://ti.arc.nasa.gov/tech/dash/groups/pcoe/prognostic-data-repository/}

\bibitem{ref10}
N. V. Chawla, K. W. Bowyer, L. O. Hall, and W. P. Kegelmeyer,
``SMOTE: Synthetic minority over-sampling technique,''
\textit{Journal of Artificial Intelligence Research}, vol.~16, pp.~321--357, 2002.

\end{thebibliography}

\end{document}